\chapter{Scroll, virtualisation, and viewports}
\label{ch:scroll}

\begin{aside}
Sooner or later your UI needs to show more content than fits the
terminal. The naive answer (render everything, clip what does
not fit) is wasteful for a list of ten thousand items.
\maya{}'s scroll model is a small data structure
(\code{ScrollState}) plus a viewport-aware rendering pattern.
This chapter walks it.
\end{aside}

\section{The ScrollState struct}

\file{maya/include/maya/core/scroll\_state.hpp} defines:

\begin{lstlisting}
struct ScrollState {
    int offset = 0;          // first visible item index
    int viewport_size = 0;   // number of rows visible
    int total_items = 0;     // total number of items in the source

    bool can_scroll_up()   const { return offset > 0; }
    bool can_scroll_down() const {
        return offset + viewport_size < total_items;
    }
    void scroll_to(int item_index, bool center = false);
    void scroll_by(int delta);
    void clamp();
};
\end{lstlisting}

Three integers (offset, viewport size, total). Every scrollable
widget stores one of these in its model. The widget's
\code{view} renders only items \code{[offset, offset+viewport\_size)}.

\subsection*{Why integers, not pointers}

A scroll position represented as ``a pointer to the visible row''
breaks the moment the underlying data changes. An integer offset
into a vector is stable: insertions and deletions can adjust it
deterministically.

\code{scroll\_to(index, center)} centres or pins the item at
\code{index}; \code{scroll\_by(delta)} moves by a relative
amount and clamps; \code{clamp()} ensures the offset is in
\code{[0, total\_items - viewport\_size]}.

\section{Virtualisation: rendering only what is visible}

A list of 10\,000 items, rendered naively, produces an element
tree with 10\,000 children. Layout walks them all. The canvas
diff walks them. ANSI emit walks them. Everything is slow.

The fix: render only the visible slice.

\begin{lstlisting}
Element render_list(const std::vector<Item>& items,
                    const ScrollState& s) {
    std::vector<Element> visible;
    for (int i = 0; i < s.viewport_size; ++i) {
        int idx = s.offset + i;
        if (idx >= s.total_items) break;
        visible.push_back(render_item(items[idx]));
    }
    return v(std::move(visible));
}
\end{lstlisting}

The element tree has at most \code{viewport\_size} children, no
matter how many items the source has. Layout, diff, emit ---
all proportional to the screen, not the data.

\begin{idea}
Virtualisation is one of the few optimisations that scales
\emph{better} than the underlying data structure. A list with a
million items renders as fast as one with twenty, because only
the twenty visible are ever materialised as elements.
\end{idea}

\subsection*{Handling tall items}

If items have variable height, the simple ``visible =
[offset, offset+viewport\_size)'' breaks: the slice's row count
depends on which items fall in it. Two approaches:

\begin{itemize}[leftmargin=*]
  \item \textbf{Constant-height items.} Easy: every item is one
    row. The slice is exactly \code{viewport\_size} items. This
    is what tables and pickers use.
  \item \textbf{Variable-height items.} Maintain a side table of
    per-item heights. To translate ``which item is on row $r$?''
    use a prefix sum or a Fenwick tree.
\end{itemize}

Most \maya{} widgets fix the height to one row. When you need
multi-row items, the widget either pre-measures or accepts the
$O($items$)$ pre-pass.

\section{Scroll bars}

The \code{scrollbar} widget reads a \code{ScrollState} and
renders a thumb:

\begin{lstlisting}
Element scrollbar(const ScrollState& s, int height) {
    if (s.total_items <= s.viewport_size) return text(\" \");
    int thumb_height = std::max(1, height * s.viewport_size / s.total_items);
    int thumb_y      = height * s.offset / s.total_items;
    // ... draw a vertical bar with thumb at thumb_y, thumb_height tall
}
\end{lstlisting}

The thumb size scales with the fraction of visible content; the
thumb position with the scroll offset. Standard math.

The scrollbar is decorative; the actual scrolling happens via
\code{ScrollState} mutations in \code{update}.

\section{Following the cursor}

A common pattern: a list with a current selection. When the
selection moves, the viewport should follow it.

\begin{lstlisting}
[&](KeyDown) {
    if (m.selected < m.items.size() - 1) m.selected++;
    if (m.selected >= m.scroll.offset + m.scroll.viewport_size)
        m.scroll.offset = m.selected - m.scroll.viewport_size + 1;
    return std::pair{m, Cmd<Msg>{}};
}
\end{lstlisting}

If the selection moves off the bottom of the viewport, scroll
down to keep it visible. The same logic applied symmetrically for
up movement keeps the cursor on-screen.

\subsection*{PageUp/PageDown}

Move the offset by \code{viewport\_size}:

\begin{lstlisting}
[&](PageDown) {
    m.scroll.scroll_by(m.scroll.viewport_size);
    m.selected = std::min(m.items.size() - 1,
                          m.selected + m.scroll.viewport_size);
    return std::pair{m, Cmd<Msg>{}};
}
\end{lstlisting}

Whether the cursor moves with the page or stays at a fixed
viewport offset is a UX choice; both are common.

\section{Horizontal scroll}

Some content (wide tables, log lines) extends past the right
edge. \code{ScrollState} can be reused on the horizontal axis:

\begin{lstlisting}
struct WideListModel {
    std::vector<std::string> lines;
    ScrollState vscroll;  // which rows are visible
    ScrollState hscroll;  // which columns
};
\end{lstlisting}

The render function clips each visible line to the horizontal
window:

\begin{lstlisting}
std::string clip_horizontal(std::string_view line, const ScrollState& h) {
    GraphemeWalker w{line};
    std::string out;
    int x = 0;
    while (auto g = w.next()) {
        if (x >= h.offset && x < h.offset + h.viewport_size) {
            out += g->text;
        }
        x += g->width;
        if (x >= h.offset + h.viewport_size) break;
    }
    return out;
}
\end{lstlisting}

Grapheme-aware, width-aware, allocation-conscious. The
performance is linear in visible cells, not in total line
width.

\section{Smooth scrolling}

For animated scroll (the viewport eases to a new position over
several frames), keep an animation state in the model:

\begin{lstlisting}
struct Scroll {
    int target = 0;       // where we want to go
    double current = 0;   // where we are now (fractional)
};
\end{lstlisting}

Each frame, advance \code{current} toward \code{target} by a
fraction of the gap. Round to integer for \code{offset}. Stop
when the difference is below 1.

The animation is driven by \code{Sub::animation\_frame}: as long
as \code{current != target}, the subscription is active.

\begin{warning}
Smooth scrolling on a terminal can look jittery because the
cell grid is integer. Snap to integer cells; smooth animations
of fractional offsets render as the same integer for multiple
frames in a row, which the user perceives as ``stuck''.
\end{warning}

\section{Sticky elements}

A header that stays visible while content scrolls under it is a
stylistic feature with one rule: layout the header separately
and never include it in the scrolled region.

\begin{lstlisting}
v(
  header() | height<1>,            // never scrolls
  scrolled_list(items, s) | grow<1>
)
\end{lstlisting}

The list grows to fill the rest of the column; the header stays
put.

For columns (a left rail that stays put while the main view
scrolls), the same pattern works on the horizontal axis.

\section{Snap-to-cell pitfalls}

\begin{warning}
\textbf{Offsets in graphemes, not bytes.} Scrolling a long line
horizontally must count graphemes, not bytes, or wide
characters will partially render.
\end{warning}

\begin{warning}
\textbf{Resize and ScrollState.} When the terminal resizes, the
viewport size changes. Always call \code{clamp()} to keep the
offset in bounds.
\end{warning}

\begin{warning}
\textbf{Dynamic total\_items.} If the source list grows or
shrinks (e.g. new log lines arrive), update
\code{total\_items} and clamp. Otherwise the user can scroll
past the end into empty space.
\end{warning}

\begin{warning}
\textbf{Offset becoming stale.} If items are inserted before
\code{offset}, the user's position drifts. Decide whether the
user's intent is ``stick with this item'' (adjust offset) or
``stay at this viewport position'' (keep offset). \maya{}'s
default is the former; sometimes you want the latter.
\end{warning}

\section{Performance numbers}

A virtualised list with a million items:

\begin{itemize}[leftmargin=*]
  \item Element tree size: bounded by viewport (e.g. 24
    children).
  \item Layout time: tens of microseconds.
  \item Diff time: sub-microsecond (only the viewport diffs).
  \item Bytes emitted per frame: a few hundred for incremental
    scroll.
\end{itemize}

The scaling is by viewport, not by source size. This is the
single optimisation that makes large data feasible in a TUI.

\section*{Workshop}

\begin{enumerate}[leftmargin=*]
  \item Build a list widget that virtualises 10\,000 items.
    Verify with a profiler that frame time is independent of
    item count.
  \item Add a scroll bar to the list using the \code{scrollbar}
    widget. Verify the thumb size and position match the
    \code{ScrollState}.
  \item Implement PageUp/PageDown that keep the selection
    centred when possible (only push it to the edge near the
    list bounds).
\end{enumerate}

\begin{recap}
\code{ScrollState} is three integers: offset, viewport size,
total. Virtualisation renders only the visible slice; layout,
diff, and emit cost are O(viewport). Sticky elements layout
outside the scrolled region. Animations are driven by signals
or \code{AnimationFrame}. The model is data; the runtime stays
unaware.
\end{recap}
